\documentclass[11pt,a4paper]{article}
\usepackage{geometry,graphicx,booktabs,longtable,caption}
\usepackage[utf8]{inputenc}
\geometry{margin=2.5cm}
\pagenumbering{arabic}

\begin{document}

\title{Supplementary Materials\\\large Nutritional Indices and All-Cause Mortality in GI Cancer: A Dual-Cohort NHANES Analysis}
\author{}
\date{}
\maketitle

\begin{center}
\textbf{Tables S1--S7, Figures S1--S3}
\end{center}
\vskip 0.5cm

% =====================================================================
\section*{Table S1: GI Cancer Site Distribution and Mortality Outcomes (Restricted to Age $\ge$60)}
\label{tab:s1}
% =====================================================================
\noindent{\footnotesize Note: This table reflects the age $\ge$60 subset (n=313, 169 deaths) from a prior analytical version. See the main text (n=353, all ages) for the primary analysis sample.}
\begin{center}
\begin{tabular}{lcccrc}
\toprule
Site & N & \% & Deaths & PNI (mean$\pm$SD) & Median survival \\
\midrule
Colon            & 175 & 55.9 & 100 & $49.8\pm4.9$ & 7.4 yr \\
Rectum           &  66 & 21.1 &  36 & $51.2\pm5.8$ & 9.2 yr \\
Liver            &  23 &  7.3 &  15 & $48.1\pm6.2$ & 5.8 yr \\
Stomach          &  19 &  6.1 &  10 & $50.5\pm5.1$ & 8.1 yr \\
Esophageal       &  15 &  4.8 &   7 & $47.9\pm5.7$ & 4.3 yr \\
Pancreatic       &   8 &  2.6 &   5 & $49.0\pm4.8$ & 6.0 yr \\
Gallbladder      &   6 &  1.9 &   3 & $51.3\pm5.5$ & 9.5 yr \\
Unknown GI       &   1 &  0.3 &  -- & --           & -- \\
\midrule
Total            & 313 & 100  & 169 & $50.8\pm6.5$ & 8.8 yr \\
\bottomrule
\end{tabular}
\end{center}

% =====================================================================
\section*{Table S2: Full Baseline Characteristics by GI Tumor Status (Restricted to Age $\ge$60)}
\label{tab:s2}
% =====================================================================
\noindent{\footnotesize Note: Restricted to age $\ge$60 (n=313 GI patients). See main text Table~1 for the full all-ages baseline.}
\begin{center}
\begin{tabular}{lccc}
\toprule
Variable & GI Tumor & Non-Cancer & Other Cancer \\
\midrule
Age, years & $67.5\pm13.8$ & $45.5\pm18.4$ & $65.6\pm15.0$ \\
Female (\%) & 64 (20.4) & 7,172 (17.7) & 1,876 (57.3) \\
Race/ethnicity (\%) & & & \\
\hspace{0.5cm}NH-White & 160 (51.1) & 14,887 (36.8) & 2,131 (65.1) \\
\hspace{0.5cm}NH-Black &  45 (14.4) & 5,909 (14.6) & 375 (11.5) \\
\hspace{0.5cm}Hispanic & 101 (32.3) & 16,394 (40.5) & 659 (20.1) \\
\hspace{0.5cm}Other &  7 (2.2) & 3,290 (8.1) & 108 (3.3) \\
Education $\ge$ HS (\%) & 187 (59.7) & 25,808 (63.8) & 2,228 (68.1) \\
BMI, kg/m$^2$ & $28.5\pm6.1$ & $28.3\pm6.6$ & $28.4\pm6.3$ \\
Albumin, g/dL & $4.1\pm0.4$ & $4.2\pm0.4$ & $4.1\pm0.4$ \\
Lymphocyte, $\mu$L$^{-1}$ & $1,972\pm990$ & $2,209\pm719$ & $1,989\pm790$ \\
Cholesterol, mg/dL & $195\pm44$ & $196\pm43$ & $199\pm43$ \\
PNI & $50.8\pm6.5$ & $53.2\pm5.2$ & $51.4\pm6.0$ \\
CONUT & $1.1\pm1.4$ & $0.7\pm1.0$ & $1.0\pm1.2$ \\
GNRI & $115.0\pm12.2$ & $116.4\pm12.4$ & $115.8\pm12.0$ \\
Deaths (\%) & 169 (54.0) & 7,659 (18.9) & 1,042 (31.8) \\
Median f/up & 8.8 yr & 13.3 yr & 10.5 yr \\
\bottomrule
\end{tabular}
\end{center}

% =====================================================================
\section*{Table S3: Baseline Characteristics by PNI Tertile Among GI Cancer Patients (Age $\ge$60 Subset)}
\label{tab:s3}
% =====================================================================
\noindent{\footnotesize Note: Restricted to age $\ge$60 (n=313).}
\begin{center}
\begin{tabular}{lccc}
\toprule
Variable & Low PNI (T1) & Mid PNI (T2) & High PNI (T3) \\
\midrule
N & 105 & 104 & 104 \\
Age, years & $72.6\pm9.4$ & $66.4\pm15.5$ & $63.6\pm13.8$ \\
Female (\%) & 15 (14.3) & 26 (25.0) & 23 (22.1) \\
Non-Hispanic White (\%) & 52 (49.5) & 57 (54.8) & 51 (49.0) \\
Education $\ge$ high school (\%) & 64 (61.0) & 58 (55.8) & 65 (62.5) \\
PNI & $44.8\pm3.5$ & $50.5\pm1.4$ & $57.5\pm4.1$ \\
CONUT & $2.1\pm1.7$ & $0.9\pm0.9$ & $0.3\pm0.5$ \\
GNRI & $110.2\pm12.6$ & $115.5\pm10.4$ & $119.5\pm11.7$ \\
Albumin, g/dL & $3.8\pm0.4$ & $4.1\pm0.3$ & $4.4\pm0.3$ \\
Lymphocyte, $\mu$L\textsuperscript{-1} & $1,376\pm389$ & $1,820\pm516$ & $2,739\pm1,049$ \\
Total cholesterol, mg/dL & $184\pm45$ & $196\pm41$ & $205\pm43$ \\
BMI, kg/m\textsuperscript{2} & $28.3\pm6.1$ & $28.5\pm5.6$ & $28.7\pm6.6$ \\
Deaths (\%) & 74 (70.5) & 48 (46.2) & 47 (45.2) \\
Median survival & 5.5 yr & 7.8 yr & 13.9 yr \\
\bottomrule
\end{tabular}
\end{center}

% =====================================================================
\section*{Table S4: Cohort-Stratified Analysis (Age $\ge$60 Subset)}
\label{tab:s4}
% =====================================================================
\noindent{\footnotesize Note: Restricted to age $\ge$60 (n=313). The all-ages analysis (n=353) showed similar non-significant interaction p-values.}
\begin{center}
\begin{tabular}{lccccc}
\toprule
Index & Cohort & N & Events & HR (95\% CI) per 1-SD & p \\
\midrule
PNI   & NHANES III & 88 & 68 & 0.69 (0.51--0.93) & 0.015 \\
PNI   & Continuous NHANES & 225 & 101 & 0.75 (0.56--1.00) & 0.047 \\
CONUT & NHANES III & 88 & 68 & 0.72 (0.56--0.93) & 0.011 \\
CONUT & Continuous NHANES & 225 & 101 & 0.69 (0.57--0.85) & $<$0.001 \\
GNRI  & NHANES III & 88 & 68 & 0.64 (0.48--0.84) & 0.001 \\
GNRI  & Continuous NHANES & 225 & 101 & 0.68 (0.55--0.84) & $<$0.001 \\
\bottomrule
\end{tabular}
\end{center}

Cohort-by-index interaction p-values: PNI p=0.480, CONUT p=0.940, GNRI p=0.950.

% =====================================================================
\section*{Table S5: Extended Covariate Sensitivity Analysis (Age $\ge$60 Subset)}
\label{tab:s5}
% =====================================================================
\noindent{\footnotesize Note: Restricted to age $\ge$60 (n=313). Results were consistent in the all-ages analysis (n=353, data not shown).}
\begin{center}
\resizebox{\textwidth}{!}{%
\begin{tabular}{lccccc}
\toprule
Model & Index & N & Events & HR (95\% CI) per 1-SD & p \\
\midrule
Basic (age+sex+race) & PNI   & 313 & 169 & 0.78 (0.63--0.96) & 0.020 \\
$+$ Education + PIR  & PNI   & 313 & 169 & 0.76 (0.62--0.94) & 0.012 \\
$+$ Smoking + Comorbidity$^\dagger$ & PNI   & 133 & 61 & 0.68 (0.48--0.97) & 0.035 \\
\midrule
Basic (age+sex+race) & CONUT & 313 & 169 & 0.71 (0.61--0.83) & $<$0.001 \\
$+$ Education + PIR  & CONUT & 313 & 169 & 0.71 (0.60--0.82) & $<$0.001 \\
$+$ Smoking + Comorbidity$^\dagger$ & CONUT & 133 & 61 & 0.63 (0.50--0.80) & $<$0.001 \\
\midrule
Basic (age+sex+race) & GNRI  & 313 & 169 & 0.67 (0.57--0.80) & $<$0.001 \\
$+$ Education + PIR  & GNRI  & 313 & 169 & 0.66 (0.56--0.79) & $<$0.001 \\
$+$ Smoking + Comorbidity$^\dagger$ & GNRI  & 133 & 61 & 0.72 (0.55--0.93) & 0.013 \\
\bottomrule
\end{tabular}}
\end{center}

\vskip3pt
\parbox{\linewidth}{\footnotesize PIR: income-to-poverty ratio. $^\dagger$Smoking and comorbidity data only available in Continuous NHANES 2005--2016 subset.}

% =====================================================================
\section*{Table S6: Survey-Weighted Cox Regression---GI Subset (Continuous NHANES 2005--2016)}
\label{tab:s6}
% =====================================================================

\begin{center}
\small
\begin{tabular}{lccccc}
\toprule
Index & N & Events & HR (95\% CI) per 1-SD & p \\
\midrule
PNI   & 270 & 134 & 0.87 (0.61--1.23) & 0.432 \\
CONUT & 270 & 134 & 0.76 (0.63--0.91) & 0.003 \\
GNRI  & 270 & 134 & 0.80 (0.61--1.05) & 0.103 \\
\bottomrule
\end{tabular}
\end{center}

See Table~3 in the main text for side-by-side comparison with unweighted results.
Survey weights: WTMEC2YR/(n cycles). Adjusted for age, sex, race/ethnicity.

% =====================================================================
\section*{Table S7: Joint PNI$\times$NLR, CRP Interaction, and CALLY Index}
\label{tab:s7}
% =====================================================================

\subsection*{S7a: Joint PNI $\times$ NLR Analysis (n=258, events=125)}

\begin{center}
\begin{tabular}{lccc}
\toprule
Group & HR (95\% CI) & p \\
\midrule
High PNI + Low NLR (ref) & 1.00 & -- \\
Low PNI + Low NLR  & 1.03 (0.61--1.75) & 0.902 \\
High PNI + High NLR & 0.79 (0.41--1.54) & 0.492 \\
Low PNI + High NLR  & 1.72 (1.06--2.63) & 0.028 \\
\bottomrule
\end{tabular}
\end{center}

RERI = $-$0.97 (95\% CI $-$2.61 to $-$0.10). The wide confidence interval limits precision; results are suggestive of sub-multiplicative interaction but should be interpreted cautiously.

\subsection*{S7b: CRP Interaction (n=250, events=120)}

\begin{center}
\begin{tabular}{lccc}
\toprule
Predictor & HR (95\% CI) & p \\
\midrule
PNI per 1-SD & 0.84 (0.63--1.10) & 0.205 \\
log(CRP) per 1-SD & 1.18 (0.92--1.50) & 0.186 \\
PNI $\times$ CRP interaction & -- & 0.780 \\
\bottomrule
\end{tabular}
\end{center}

\subsection*{S7c: CALLY Index Comparison (n=250, events=120)}

\begin{center}
\begin{tabular}{lcccc}
\toprule
Index & HR (95\% CI) & p & C-statistic \\
\midrule
PNI & 0.85 (0.67--1.09) & 0.200 & 0.667 \\
CALLY & 0.85 (0.67--1.09) & 0.200 & 0.659 \\
\bottomrule
\end{tabular}
\end{center}

CALLY = albumin $\times$ lymphocyte / CRP. Neither index reached statistical significance in this subset with CRP data. C-statistic difference was negligible ($\Delta$=0.008), indicating no advantage of CALLY over PNI.

% =====================================================================
\section*{Table S8: GI Cancer Patient Characteristics by NHANES Generation}
\label{tab:s8}
% =====================================================================

\begin{center}
\small
\begin{tabular}{lccc}
\toprule
Variable & NHANES III (1988--1994, n=126) & Continuous NHANES (2005--2016, n=227) & p \\
\midrule
Age, years & 67.5 (13.9) & 57.5 (19.0) & $<$0.001 \\
Female (\%) & 43.1 & 59.7 & 0.002 \\
Non-Hispanic White (\%) & 68.4 & 32.6 & $<$0.001 \\
BMI, kg/m$^2$ & 29.3 (6.1) & 27.7 (6.5) & 0.014 \\
Albumin, g/dL & 4.12 (0.36) & 4.06 (0.41) & 0.178 \\
PNI & 50.7 (6.5) & 51.7 (6.0) & 0.113 \\
CONUT & 1.28 (1.27) & 0.69 (1.25) & $<$0.001 \\
GNRI & 116.8 (11.7) & 113.0 (12.3) & 0.002 \\
Site distribution (\%) & & & $<$0.001 \\
\quad Colorectal & 40.5 & 79.3 \\
\quad Liver & 15.1 & 4.4 \\
\quad Upper GI & 14.3 & 11.9 \\
\quad Pancreatic & 0.0 & 3.1 \\
\quad Unknown & 30.2 & 1.3 \\
Deaths (\%) & 78.6 & 48.0 & $<$0.001 \\
Median survival, yr & 19.1 [7.0--26.8] & 7.1 [4.3--10.8] & $<$0.001 \\
Early death ($<$2 yr, \%) & 11.7 & 5.5 & 0.042 \\
\bottomrule
\end{tabular}
\end{center}

Values are mean (SD) unless otherwise indicated. The markedly longer median follow-up in NHANES III reflects the longer elapsed time between survey participation (1988--1994) and mortality ascertainment (2019). No statistically significant heterogeneity was detected across generations (all cohort-by-index interaction p$>$0.05; Table~S4), though the modest sample sizes limit the power of these interaction tests.

% =====================================================================
\section*{Supplementary Results: Dietary Quality (HEI-2015)}
\label{sec:hei}
% =====================================================================

Among 99 GI cancer patients with available dietary recall data (Continuous NHANES cycles with Food Patterns Equivalents Database), HEI-2015 did not predict all-cause mortality (HR 0.97 per 1-SD, 95\% CI 0.70--1.35, p=0.850). The correlation between HEI-2015 and PNI was weak (r=0.16), confirming that biochemical nutritional indices and dietary intake patterns capture distinct dimensions of nutritional status. However, these results are limited by a single 24-hour dietary recall (which does not reflect usual intake), the small sample size (n=99, 58 events), and potential selection bias in the subset with complete dietary data. The null finding should therefore be interpreted cautiously as it is equally consistent with a true null effect and with insufficient statistical power.

\end{document}
